\subsection{The Architecture of Non-Local Jumps}

Ordinary control flow moves locally. A sequential statement falls into the next statement. A condition chooses one nearby branch. A loop returns to its own condition or iterator. An exception is different. When an exception is raised, execution may leave the current expression, the current statement, the current block, and even the current function call.

This is why exceptions are a form of non-local control transfer. The next relevant statement may not be physically near the failing operation. Python begins searching for a matching exception handler. If none exists in the current block, the exception can continue outward through active function calls until it is caught or until it reaches the interpreter boundary.

\begin{quote}
An exception is not merely an error value. It is a runtime event that changes the execution route.
\end{quote}

This section introduces exceptions as routing structures. Later subsections will examine exception classes, handler matching, \texttt{try}/\texttt{except}/\texttt{else}/\texttt{finally}, traceback construction, explicit \texttt{raise}, re-raising, and exception chaining.

\subsubsection{The Paradigm Shift: Contrast with C's Manual Error Return Codes and Pointer Check Patterns}

In C, failure is often represented through returned values or modified output parameters. A function may return \texttt{-1}, \texttt{NULL}, \texttt{0}, or an error code. The caller must inspect that result manually before continuing.

A simplified C-style pattern looks like this:

%<<<PROGRAM:programs/the11>>>
The failure information stays local unless the programmer explicitly forwards it. Every caller must remember to check the returned value. If the programmer forgets a check, the program may continue with invalid state.

Another common C-style pattern is pointer checking:

\begin{verbatim}
/* C-style idea, not Python */
char *name = get_name();

if (name != NULL) {
    printf("%s\n", name);
}
\end{verbatim}

The programmer manually asks whether the pointer is usable before dereferencing it. The safety of the next operation depends on the explicit check.

Python can also use ordinary conditional checks, but its exception mechanism changes the default failure route. If an operation cannot complete normally, it can raise an exception instead of returning a special error value.

%<<<PROGRAM:programs/the10>>>
The expression \texttt{int(text)} does not return a special error number. It raises a \texttt{ValueError}. The ordinary route inside the \texttt{try} block is abandoned, and execution jumps to the matching \texttt{except} block.

The exception object is a runtime object:

%<<<PROGRAM:programs/the09>>>
The \texttt{args} attribute stores the construction arguments carried by the exception. The \texttt{\_\_traceback\_\_} attribute connects the exception to traceback information. Later sections will examine how tracebacks preserve the route of failure through active frames.

The important control-flow difference is this:

\begin{verbatim}
C-style error code:
    operation returns an error marker
    caller manually checks marker
    caller manually chooses next route

Python exception:
    operation raises an exception
    ordinary route is interrupted
    Python searches for a matching handler
\end{verbatim}

This does not mean that Python eliminates responsibility. The programmer still decides which exceptions to catch, where to catch them, and what recovery means. The difference is that failure travels through a separate control-flow channel rather than through an ordinary returned value.

A simple dictionary example shows the same idea:

\begin{verbatim}
scores = {
    "Jerry": 72,
    "Elaine": 91,
    "Kramer": 42,
}

name = "Newman"

try:
    score = scores[name]
    print(f"{name}: {score}")
except KeyError as err:
    print(f"missing key: {err}")
    print(f"type(err): {type(err)}")

print("lookup attempt finished")
\end{verbatim}

Possible output:

\begin{verbatim}
missing key: 'Newman'
type(err): <class 'KeyError'>
lookup attempt finished
\end{verbatim}

The dictionary object is an ordinary runtime object:

%<<<PROGRAM:programs/the08>>>
The \texttt{\_\_getitem\_\_()} method is involved in bracket lookup such as \texttt{scores[name]}. If the key is absent, that lookup raises \texttt{KeyError}. The \texttt{get()} method offers a different design: return a fallback value instead of raising.

%<<<PROGRAM:programs/the07>>>
Both approaches are valid in different situations. The bracket form treats a missing key as an exceptional route. The \texttt{get()} form treats a missing key as an ordinary value-selection route.

\subsubsection{Look Before You Leap (LBYL) vs. Easier to Ask Forgiveness than Permission (EAFP) Execution Philosophies}

Two common styles appear in Python error-handling discussions: LBYL and EAFP.

LBYL means ``look before you leap.'' In this style, the program checks whether an operation appears safe before performing it.

%<<<PROGRAM:programs/the06>>>
The condition checks membership before the lookup. The lookup happens only if the key is present.

EAFP means ``easier to ask forgiveness than permission.'' In this style, the program attempts the operation directly and handles the exception if the operation fails.

\begin{verbatim}
scores = {
    "Jerry": 72,
    "Elaine": 91,
    "Kramer": 42,
}

name = "Newman"

try:
    score = scores[name]
    print(f"{name}: {score}")
except KeyError:
    print(f"{name} was not found")
\end{verbatim}

Possible output:

\begin{verbatim}
Newman was not found
\end{verbatim}

Both programs produce the same visible result. Their control-flow shape is different.

The LBYL route is:

\begin{verbatim}
test whether key exists
    true  -> perform lookup
    false -> run missing-key branch
\end{verbatim}

The EAFP route is:

\begin{verbatim}
attempt lookup
    success  -> continue normally
    KeyError -> jump to handler
\end{verbatim}

The EAFP version avoids duplicating the key lookup logic. This can be useful when the safest test is the operation itself.

A list example makes the difference clear:

%<<<PROGRAM:programs/the05>>>
That is the LBYL version. It checks the index before using it.

The EAFP version is:

%<<<PROGRAM:programs/the04>>>
The list object supports indexed access through \texttt{\_\_getitem\_\_()}:

\begin{verbatim}
items = ["soup", "pretzels"]

print(f"type(items): {type(items)}")
print(f"dir(items) sample: {dir(items)[:10]}")
print(f"len(items): {len(items)}")
print(f"items.__getitem__(0): {items.__getitem__(0)}")
\end{verbatim}

Possible output:

\begin{verbatim}
type(items): <class 'list'>
dir(items) sample: ['__add__', '__class__', '__class_getitem__',
'__contains__', '__delattr__', '__delitem__', '__dir__', '__doc__',
'__eq__', '__format__']
len(items): 2
items.__getitem__(0): soup
\end{verbatim}

If the requested index is invalid, the indexed operation raises \texttt{IndexError}. The exception route is part of the object's interface behavior.

Neither LBYL nor EAFP is universally correct. They are different ways to arrange control flow.

LBYL is often clearer when the condition is cheap, direct, and part of the intended ordinary logic:

%<<<PROGRAM:programs/the03>>>
Here the empty-string check expresses an ordinary input condition. There is no need to force that case through an exception.

EAFP is often clearer when the operation itself is the authoritative test:

%<<<PROGRAM:programs/the02>>>
Checking every possible invalid integer string manually would be awkward and incomplete. The conversion operation already knows the valid format. Letting \texttt{int()} decide and catching \texttt{ValueError} gives a precise route.

A useful combined version is:

%<<<PROGRAM:programs/the01>>>
The empty-input case is ordinary branching. The invalid-format case is exceptional routing.

The practical rule is:

\begin{quote}
Use LBYL when the precondition is simple and genuinely part of normal decision logic. Use EAFP when the attempted operation is the most accurate way to discover whether it can succeed.
\end{quote}

For control-flow graph analysis, the distinction is structural. LBYL adds explicit conditional edges before the operation. EAFP adds exceptional edges after the operation begins. Both are routing tools. The correct choice depends on which route makes the program's intended facts clearer.